% cap05_clifford_rotores.tex
\chapter{Álgebra de Clifford y Rotores Geométricos}


El estudio de rotaciones y transformaciones ortogonales en espacios euclídeos de alta dimensión, $\mathbb{R}^D$, requiere de herramientas algebraicas robustas y computacionalmente eficientes. Mientras que la representación tradicional mediante matrices ortogonales de grupo $O(D)$ ha dominado la literatura debido a su inmediata implementación en hardware lineal, presenta severas limitaciones en entornos hiperdimensionales ($D \ge 10,000$). La Álgebra de Clifford proporciona un marco unificado que no solo supera estos escollos asintóticos, sino que también revela la estructura topológica subyacente de $S^{D-1}$, siendo la base fundamental de la Programación Cognitiva en POLYDIM.

\section{Definición del Álgebra de Clifford $Cl(D,0)$}

El álgebra de Clifford $Cl(D,0)$ sobre un espacio vectorial $V$ de dimensión $D$ con una forma cuadrática definida positiva $Q(x) = \|x\|^2$ es el álgebra asociativa con unidad generada por los elementos de $V$, sujeta a la relación fundamental:
\begin{equation}
    x^2 = Q(x)1 = \|x\|^2 \quad \forall x \in V
\end{equation}
De esta identidad, mediante el proceso de polarización con $x, y \in V$, se deduce inmediatamente que:
\begin{equation}
    (x+y)^2 = x^2 + xy + yx + y^2 = Q(x+y)
\end{equation}
Dado que $Q(x+y) = \|x+y\|^2 = \|x\|^2 + \|y\|^2 + 2\langle x, y \rangle$, obtenemos la relación de anticonmutación:
\begin{equation}
    xy + yx = 2\langle x, y \rangle
\end{equation}

Esta simple regla gobierna toda la estructura geométrica. En particular, si $x$ e $y$ son ortogonales ($\langle x, y \rangle = 0$), entonces anticomutan: $xy = -yx$.

\section{El Producto Geométrico}

El producto en $Cl(D,0)$ se denomina producto geométrico. Para dos vectores $a,b \in V$, este producto se puede descomponer unívocamente en una parte simétrica y una parte antisimétrica:
\begin{equation}
    ab = \frac{1}{2}(ab + ba) + \frac{1}{2}(ab - ba)
\end{equation}
Reconociendo que la parte simétrica coincide con el producto interno (escalar), y definiendo la parte antisimétrica como el producto exterior (producto cuña o \textit{wedge product}), obtenemos la identidad central del álgebra geométrica:
\begin{equation}
    ab = a \cdot b + a \wedge b
\end{equation}
Donde:
\begin{itemize}
    \item $a \cdot b = \langle a, b \rangle$ es un escalar (grado 0), que mide la colinealidad.
    \item $a \wedge b$ es un bivector (grado 2), que representa el segmento de plano orientado definido por $a$ y $b$.
\end{itemize}
Esta dualidad simultánea es la que dota al producto geométrico de su inmenso poder: codifica tanto métrica (producto interno) como subespacios y orientaciones (producto exterior) sin recurrir a tensores de orden superior de forma explícita.

\section{Construcción a través del Álgebra Exterior y Estructura de Grados}

Formalmente, como espacio vectorial, el álgebra de Clifford $Cl(D,0)$ es isomorfa al álgebra exterior $\Lambda(V)$. Su dimensión total es $2^D$, y se descompone en suma directa de subespacios de grado $k$ (o \textit{k-blades}):
\begin{equation}
    Cl(D,0) = \bigoplus_{k=0}^{D} \Lambda^k(V)
\end{equation}
La base canónica para un $k$-vector se forma a partir del producto exterior de $k$ vectores ortonormales de la base de $V$: $\{e_1, e_2, \dots, e_D\}$. Un multivector genérico $M$ se expresa como:
\begin{equation}
    M = \langle M \rangle_0 + \langle M \rangle_1 + \langle M \rangle_2 + \dots + \langle M \rangle_D
\end{equation}
Donde $\langle M \rangle_0$ es la parte escalar (el invariante de grado 0), de crucial importancia para el protocolo de consenso distribuido en el Swarm. En POLYDIM, evitamos instanciar la base completa $O(2^D)$, concentrándonos en $k$-blades puros (específicamente bivectores, $k=2$) que rigen las rotaciones sobre $S^{D-1}$.

\section{Rotores: $R = \exp(-B\theta/2)$ y la Acción Bilátera}

Las rotaciones en un espacio vectorial pueden modelarse como la composición de dos reflexiones a lo largo de hiperplanos ortogonales a vectores unitarios $m$ y $n$. En $Cl(D,0)$, la reflexión de un vector $x$ respecto al hiperplano ortogonal a $n$ (con $\|n\|=1$) está dada por:
\begin{equation}
    x' = -n x n
\end{equation}
Aplicando dos reflexiones sucesivas mediante $n$ y luego $m$:
\begin{equation}
    x'' = -m (-n x n) m = (mn) x (nm)
\end{equation}
Definimos el Rotor $R = mn$. Nótese que, al ser producto de vectores, $R$ pertenece a la subálgebra par $Cl^+(D,0)$, conocida como el grupo Spin, $Spin(D)$. La acción del rotor sobre el vector $x$ es una acción bilátera (\textit{double-sided action}):
\begin{equation}
    x'' = R x \tilde{R}
\end{equation}
Donde $\tilde{R} = nm$ es el reverso (inversión del orden de los factores) de $R$. 

Por el teorema de Cartan-Dieudonné, cualquier rotación en $D$ dimensiones puede expresarse como producto de a lo sumo $D$ reflexiones. Así, cualquier matriz ortogonal especial $M \in SO(D)$ tiene su equivalente en $Spin(D)$. Además, como el plano de rotación está definido por el bivector $B = n \wedge m$, podemos expresar el rotor continuo mediante la forma exponencial:
\begin{equation}
    R = \exp\left(-\frac{\theta}{2} B\right) = \cos\left(\frac{\theta}{2}\right) - B \sin\left(\frac{\theta}{2}\right)
\end{equation}
Donde $B^2 = -1$. Esta expresión es idéntica a la fórmula de Euler, pero generalizada a bivectores en dimensiones arbitrarias.

\section{Comparativa Topológica y Computacional: $O(D)$ frente a Rotores Clifford}

\subsection{Doble Recubrimiento Universal (\textit{Double Cover})}
El grupo $Spin(D)$ es el doble recubrimiento (universal para $D \ge 3$) del grupo ortogonal especial $SO(D)$. Existe un epimorfismo $\rho: Spin(D) \to SO(D)$ con núcleo $\{-1, 1\}$. Esto significa que los rotores $R$ y $-R$ corresponden a la misma rotación en $SO(D)$. Esta propiedad métrica es la que previene problemas topológicos como el \textit{Gimbal Lock} que acosan a los ángulos de Euler, y permite interpolaciones suaves a lo largo de geodésicas en la hiperesfera (SLERP).

\subsection{Complejidad de Memoria y Cómputo}
En matrices, representar una rotación en $\mathbb{R}^D$ requiere una matriz de tamaño $D \times D$, es decir, una huella de memoria $O(D^2)$. En arquitecturas como Cerebras WSE o en los nodos del Swarm procesando tensores con $D=10^6$, almacenar $O(D^2)$ flotantes es inviable (requeriría $\sim 4$ Terabytes solo para la matriz de rotación).

Por el contrario, un rotor simple (rotación en un único plano) requiere únicamente identificar el bivector generador $B$ (definido por dos vectores, $u$ y $v$) y el ángulo $\theta$. Esto requiere $O(D)$ memoria para almacenar $u$ y $v$, colapsando la complejidad espacial. Esta es la revolución asintótica de POLYDIM: la Fórmula de Rodrigues generalizada.

\section{Conexión con POLYDIM: Generalización de la Fórmula de Rodrigues}
El núcleo en C++ `kernel_cpp_v762.cpp` y el kernel Triton implementan la acción del rotor $x' = R x \tilde{R}$ sin instanciar la estructura completa $2^D$ de multivectores, utilizando una proyección geométrica puramente vectorial. La acción bilátera se desglosa y optimiza directamente como:
\begin{equation}
    x_{rot} = x + u(\dots) + v(\dots)
\end{equation}
(detallado exhaustivamente en el Capítulo 6). De este modo, POLYDIM aprovecha el isomorfismo métrico de Clifford para lograr eficiencia $O(D)$ en cómputo y memoria, manteniendo al 100\% la precisión topológica en $S^{D-1}$.

\section{El Operador Duality Estrella de Hodge $\star$}
El elemento pseudoescalar unitario $I \in Cl(D,0)$ se define como el producto de todos los vectores de la base ortonormal: $I = e_1 e_2 \dots e_D$. Este elemento conmuta o anticonmuta con todos los demás dependiendo de la paridad de $D$. El operador de dualidad de Hodge $\star$ se relaciona íntimamente con la multiplicación por el pseudoescalar $I$. Para cualquier multivector $A$, su dual es $A^* = A I^{-1}$. En POLYDIM, esta dualidad permite mapear problemas de alta dimensión a subespacios de menor grado para ataques Red Team.

\section{Isometría Absoluta: $\|R x \tilde{R}\| = \|x\|$}
\begin{theorem}[Isometría Geométrica]
La acción bilátera de un rotor $R$ sobre un vector $x$, denotada $x' = R x \tilde{R}$, preserva exactamente la norma geométrica de $x$.
\end{theorem}
\begin{proof}
Consideremos la norma al cuadrado de $x'$, que en álgebra geométrica es el producto geométrico por su reverso. Como $x$ es vector, su reverso es sí mismo $\tilde{x} = x$, pero debemos considerar el reverso de todo el producto:
\begin{align*}
    \|x'\|^2 &= (R x \tilde{R}) (R x \tilde{R})\tilde{\ } \\
             &= (R x \tilde{R}) (R \tilde{x} \tilde{R}) \\
             &= R x (\tilde{R} R) x \tilde{R}
\end{align*}
Para un rotor, $\tilde{R} R = 1$. Por tanto:
\begin{align*}
    \|x'\|^2 &= R x (1) x \tilde{R} \\
             &= R x^2 \tilde{R}
\end{align*}
Dado que $x$ es un vector, $x^2 = \|x\|^2$, que es un escalar. Los escalares conmutan con cualquier multivector:
\begin{align*}
    \|x'\|^2 &= \|x\|^2 R \tilde{R} = \|x\|^2 (1) = \|x\|^2
\end{align*}
En consecuencia, $\|x'\| = \|x\|$. Esta prueba algebraica no depende de $D$ ni de la representación matricial matricial.
\end{proof}
Esta propiedad se verifica asintóticamente en POLYDIM, asegurando Deriva Cero (Drift = 0.0) a precisiones de máquina.

\section{El Conjunto de Compuertas Clifford+T y la Convergencia Cuántica}
El interés de formalizar los operadores en $S^{D-1}$ como elementos de $Cl(D,0)$ no es meramente abstracto. En el contexto de la computación cuántica, el teorema de Solovay-Kitaev demuestra que el conjunto de compuertas de Clifford, complementado con la compuerta $T$ (rotación de $\pi/8$), es universal.

POLYDIM anticipa la migración del cómputo de tensores a hardware cuántico y QPU's. Al enmarcar las transformaciones de estado interno (LatentMAS PMTP) puramente como rotores espaciales sobre geodésicas (sin la contaminación 1D de un Transformer), los estados de red son isomórficos a estados cuánticos puramente unitarios. El "Ghost Protocol" de transferencia multihop sin colapso a 1D (texto/JSON) preserva el vector de estado inalterado, como un sistema cuántico no observado, colapsando solo en la interfaz final "2D Worm".

\section{Consideraciones Físicas y de Hardware: FtzDazGuard RAII}
En la ejecución a nivel de silicio, la preservación isométrica $S^{D-1}$ descrita teóricamente es acosada por las limitaciones del formato IEEE-754. Las interacciones con números subnormales pueden desencadenar trampas del hardware (\textit{microcode traps}), demorando las canalizaciones (pipelines) aritméticas masivamente.
\begin{verbatim}
class FtzDazGuard {
    unsigned int original_mxcsr;
public:
    FtzDazGuard() {
        original_mxcsr = _mm_getcsr();
        // Set Flush-To-Zero (bit 15) and Denormals-Are-Zero (bit 6)
        _mm_setcsr(original_mxcsr | (1 << 15) | (1 << 6));
    }
    ~FtzDazGuard() {
        _mm_setcsr(original_mxcsr);
    }
};
\end{verbatim}
En la implementación asintótica descrita en `kernel_cpp_v762.cpp`, los rotores de Clifford demandan miles de sumas concurrentes por OpenMP. Para prevenir las anomalías de los estados subnormales (flush-to-zero y denormals-are-zero), se instila el `FtzDazGuard`, garantizando que la isometría probabilística no descarrile el ancho de banda por \textit{branching} subnormal de bajo nivel.

\section{Conclusión}
El paso del marco matricial euclídeo hacia el álgebra de Clifford purifica el cómputo de rotaciones hiperdimensionales, proveyendo un mapeo exacto $O(D)$ en lugar del catastrófico costo de memoria $O(D^2)$. El rotor como un invariante topológico de doble recubrimiento habilita explícitamente a POLYDIM para ejecutar SLERP en espacios $D \ge 10,000$ con garantías estrictas de Cero Deriva.
